% ============================================================================
%  CHAPTER 3: METHODOLOGY   (expanded, deep-dive version with pseudocode)
%  Replaces the entire Chapter 3 body in main.tex.
%  Pseudocode uses the listings package (style "pseudo") for line numbering;
%  each block is preceded by \algname{...} which prints its name in bold.
%  Keep code lines <= ~72 chars to stay inside the 1in margins.
% ============================================================================

\chapter{Methodology}
\label{chap:methodology}

This chapter describes how FinITR-AI is actually built, in enough detail that
the design could be reconstructed from the text. One principle runs through
everything that follows, and it is worth stating at the outset because it
explains almost every decision below: the language model is permitted to reason,
classify, and explain, but it is never permitted to compute a number that
matters. Every quantity with a financial consequence --- tax, surcharge, cess,
rebate, HRA exemption, the risk score --- is produced by deterministic code that
the model cannot override. Around that principle we arrange four specialised
agents, a small set of shared tools, a layer of document parsers, and two
machine-learning models. We begin by restating the problem in operational terms,
then give the overall approach and the shared data model, and then walk through
the architecture and each of the four agent modules with their algorithms set
out explicitly.

\section{Problem Restatement}
\label{sec:problem-restatement}

Stripped of narrative, the system is a function. Its input is some combination
of a taxpayer's Form 16 (a PDF or a JSON), their Annual Information Statement (a
JSON), a bank statement (a CSV), and optionally Form 26AS. Its output is a
structured, filing-ready report. The transformation in between must satisfy
three requirements, and the rest of the chapter is essentially an account of how
each is met.

\begin{itemize}
    \item It must reconcile three independent and partially overlapping views of
          the same financial year --- what the employer reported in Form 16,
          what the bank recorded transaction by transaction, and what the
          government already holds in the AIS --- and flag every material
          discrepancy as a quantified notice risk rather than a buried data
          point.
    \item It must compute the exact tax liability under both the Old and the New
          regimes for FY 2025-26, including the standard deduction, the Section
          87A rebate and its marginal-relief zone, surcharge with its own
          marginal relief, and the four percent cess, and then recommend the
          regime and the deductions that lawfully minimise what the user owes.
    \item It must select the correct ITR form, map each income source to the
          schedule it belongs in, and certify that every claim appearing in the
          final report is grounded either in a source document or in a statutory
          provision --- refusing to assert anything it cannot support.
\end{itemize}

\section{Overall Approach}
\label{sec:overall-approach}

The system is layered. At the bottom is a parsing layer that turns
heterogeneous documents into one normalised, shared object. In the middle is the
agentic reasoning core, where the intelligence lives. At the top is an output
layer that renders the enriched object into a report, an ITR-2 JSON, and a PDF
brief. The middle layer is itself organised around three ideas that recur
throughout the chapter.

\begin{itemize}
    \item \textbf{Specialised agents over a shared context.} Four agents ---
          Auditor, Optimizer, Compliance, and Critic --- read from and write to a
          single shared state object called the \texttt{AgentContext}. Each
          agent builds on the verified output of the ones before it instead of
          re-deriving the world from scratch, which keeps the reasoning of each
          stage focused and inspectable.
    \item \textbf{Tools the agents must defer to.} A CalculatorTool performs all
          arithmetic, a PageIndex retriever supplies legal citations, a vector
          store enables semantic document search, and a faithfulness verifier
          checks claims. Agents decide \emph{what} to ask; the tools determine
          the answer. An agent cannot, by construction, invent a tax figure.
    \item \textbf{A ReAct control loop with a critic inside it.} The orchestrator
          runs the working agents in sequence and then lets the Critic review the
          assembled result. If a claim cannot be grounded, the orchestrator
          re-runs only the responsible agent --- not the whole pipeline --- with
          the critic's constraints attached, for a bounded number of iterations.
\end{itemize}

\section{System Block Diagram and Data Model}
\label{sec:block-diagram}

The block diagram in Figure~\ref{fig:block-diagram} shows the high-level
structure and the flow of data. Raw documents enter on the left; parsers convert
them into the shared context; the orchestrated agents enrich that context; two
machine-learning models annotate it; and the output layer renders the final
artefacts.

\begin{figure}[H]
    \centering
    \includegraphics[width=0.9\textwidth]{images/block_diagram.png}
    \caption{System block diagram --- document ingestion, the shared AgentContext, the four-agent core, the inline ML models, and the output layer.}
    \label{fig:block-diagram}
\end{figure}

The pivot of the whole design is the shared data model, the \texttt{AgentContext}.
It is a single mutable record that every component reads from and writes to, and
it is what allows the agents to cooperate without passing tangled arguments
around. Its principal fields are the following:

\begin{itemize}
    \item the parsed document inputs --- \texttt{bank\_transactions},
          \texttt{form16\_data}, \texttt{ais\_data}, \texttt{form26as\_data};
    \item the accumulated agent outputs --- \texttt{reconciliation} and
          \texttt{anomalies} and \texttt{risk\_score} from the Auditor,
          \texttt{regime\_comparison} and \texttt{ctc\_strategy} from the
          Optimizer, \texttt{itr\_form\_recommendation} and
          \texttt{schedule\_mapping} from the Compliance agent, and
          \texttt{verification\_results} from the Critic;
    \item the user profile --- \texttt{gross\_income}, \texttt{basic\_salary},
          \texttt{employer\_nps}, and the \texttt{assessment\_year}; and
    \item the orchestrator's own bookkeeping --- the current \texttt{iteration},
          the \texttt{max\_iterations} bound, the running \texttt{agent\_trace},
          and the \texttt{critic\_feedback} list that carries blocked claims
          between iterations.
\end{itemize}

A small but important piece of robustness lives at the parsing boundary. Real
documents are noisy: a Form 16 produced by optical character recognition will
often disagree with the AIS by a few rupees because of rounding. The
orchestrator absorbs this rather than propagating it. When the Form 16 gross
salary and the AIS salary differ by no more than three percent, it treats the
difference as noise and adopts the AIS figure as ground truth, scaling the basic
salary proportionally. The logic is shown below.

\algname{reconcile\_gross\_income(form16, ais)}
\begin{lstlisting}[style=pseudo]
PROCEDURE reconcile_gross_income(form16, ais):
    f16_gross <- form16.gross_salary
    ais_salary <- sum(e.amount for e in ais.sft_entries
                                 if e.type == "salary")
    IF ais_salary > 0 AND
       |f16_gross - ais_salary| <= 0.03 * ais_salary THEN
        gross_income <- ais_salary          # trust AIS
        IF form16.basic > 0 AND f16_gross > 0 THEN
            basic <- round(form16.basic * (ais_salary / f16_gross))
        ELSE
            basic <- 0.40 * ais_salary
    ELSE
        gross_income <- f16_gross            # genuine mismatch: keep both
        basic <- form16.basic OR 0.40 * f16_gross
    RETURN gross_income, basic
\end{lstlisting}

The decision to trust the AIS within the tolerance, rather than the
employer-issued Form 16, follows from the asymmetry described in
Chapter~\ref{chap:introduction}: the department's record is the one against
which the return is matched, so where the two nearly agree, aligning to the AIS
removes a spurious mismatch before it can inflate the risk score.

\section{Architecture and the ReAct Control Loop}
\label{sec:architecture}

The architecture, shown in Figure~\ref{fig:architecture}, is driven by the
orchestrator. Every agent is a subclass of a common base that provides timing,
tool access, and a uniform result structure, so the orchestrator can treat them
interchangeably and record, for each one, its status, its duration, the tools it
called, and any warnings it raised.

\begin{figure}[H]
    \centering
    \includegraphics[width=0.9\textwidth]{images/architecture.png}
    \caption{System architecture --- the ReAct orchestrator, the four agents, the shared deterministic tools, and the critic re-run loop.}
    \label{fig:architecture}
\end{figure}

The control loop itself is intentionally simple, and its bounded nature is a
deliberate safeguard against the kind of runaway self-correction that
multi-agent systems can fall into. On the first pass all three working agents
run and the Critic reviews. If the Critic is satisfied, the report is finalised.
If not, the Critic's feedback names the offending agent and the specific claims
it blocked; the orchestrator re-runs only that agent, with the constraints
attached to the context, and then re-runs the Critic. This continues for at most
three iterations. The full procedure is given below.

\algname{orchestrate(documents) \quad / \quad should\_rerun(agent\_name, ctx)}
\begin{lstlisting}[style=pseudo]
PROCEDURE orchestrate(documents):
    ctx <- new AgentContext
    parse_documents_into(ctx, documents)        # fills ctx inputs
    reconcile_gross_income(ctx)                  # noise tolerance

    iteration <- 0
    WHILE iteration < ctx.max_iterations:
        ctx.iteration <- iteration

        IF iteration == 0 OR should_rerun("auditor", ctx):
            run(Auditor, ctx)                    # writes reconciliation,
                                                 # anomalies, risk_score
        IF iteration == 0 OR should_rerun("optimizer", ctx):
            run(Optimizer, ctx)                  # writes regime_comparison
        IF iteration == 0 OR should_rerun("compliance", ctx):
            run(Compliance, ctx)                 # writes itr_form, schedules

        result <- run(Critic, ctx)               # verifies all claims
        IF result.status == "success":
            BREAK                                # nothing to fix
        ctx.critic_feedback.append({
            issues: result.warnings,
            blocked_claims: result.blocked_claims
        })
        iteration <- iteration + 1

    report <- build_report(ctx)
    report.notice_prediction <- NoticePredictor.predict(report)
    RETURN report

PROCEDURE should_rerun(agent_name, ctx):
    # re-run an agent only if the latest critic feedback
    # names it as the source of an issue or a blocked claim
    latest <- ctx.critic_feedback.last
    RETURN any(agent_name in issue.source
               for issue in latest.issues)
        OR any(agent_name in claim.source
               for claim in latest.blocked_claims)
\end{lstlisting}

After the loop terminates, two machine-learning models annotate the report. The
transaction classifier has already done its work inside the Auditor; the
standalone notice predictor now estimates the probability of a Section 143(1)
notice from features of the finished report. Both are described in the module
sections that follow.

\subsection{A Worked Example}
It helps to fix ideas with a concrete trace before dissecting each agent.
Consider a salaried filer with a gross salary of about Rs.~18,50,000 declared on
Form 16, whose AIS additionally contains a fixed-deposit interest entry of
Rs.~42,000 (reported under SFT but absent from Form 16) and an equity-sale
proceed of Rs.~1,20,000 routed through a broker, with a matching credit visible
in the bank statement. The pipeline proceeds as follows. The Auditor classifies
the bank credits, finds the FD interest and the equity sale in the AIS but not
in Form 16, records them in the reconciliation ledger as \texttt{ais\_only}
items, and scores the resulting notice risk. The Optimizer computes the tax both
ways; under the New Regime the standard deduction of Rs.~75,000 brings taxable
income to Rs.~17,75,000, and the slab computation (described in
Section~\ref{sec:module-two}) yields a slab tax near Rs.~1,55,000 to which the
four percent cess adds roughly Rs.~6,200, for a liability of about
Rs.~1,61,200; the Old Regime, lacking enough deductions to overcome its higher
rates, comes out higher, so the New Regime is recommended. The Compliance agent
sees a capital-gains signal and therefore selects ITR-2, mapping the equity sale
to Schedule CG and the interest to Schedule OS. The Critic verifies that no
disallowed deduction was claimed under the New Regime and that each cited section
exists, and the notice predictor, seeing AIS income not present in Form 16,
returns an elevated probability. We return to the individual steps below.

\section{Module One -- Auditor Agent}
\label{sec:module-one}

The Auditor is the entry point and, in many ways, the core contribution of the
system. Its responsibility is reconciliation and risk. Its overall workflow,
shown in Figure~\ref{fig:module1}, has six steps: extract income items from each
of the three sources, merge them into a unified ledger, cross-reference the AIS
against Form 16, detect anomalies in the bank statement, compute a risk score
from the evidence, and generate targeted interview questions for the items that
need human input.

\begin{figure}[H]
    \centering
    \includegraphics[width=0.9\textwidth]{images/module1_flow.png}
    \caption{Auditor Agent --- transaction classification, AIS-versus-Form-16 reconciliation, and additive risk scoring.}
    \label{fig:module1}
\end{figure}

\subsection{The Three-Stage Transaction Classifier}
Before reconciliation can begin, every bank credit must be assigned an income
type. This is harder than it sounds, because real Indian bank statements are
written in a terse, inconsistent shorthand --- transfer codes, UTR references,
truncated merchant names, and Hinglish vendors such as ``Aman Juicewala'' or
``Sharma Sweet Mart''. The classifier addresses this with a three-stage cascade
that trades speed for coverage: cheap rules first, learned similarity next, and
the language model only as a last resort. The taxonomy has twelve labels, each
carrying a tax-relevance flag and an ITR schedule.

\algname{classify\_transaction(description, direction)}
\begin{lstlisting}[style=pseudo]
PROCEDURE classify_transaction(description, direction):
    cleaned <- normalize(description, direction)   # strip codes/refs

    # Stage 1 -- high-precision regex (resolves the majority,
    #            sub-millisecond, no model loaded)
    FOR (pattern, label, confidence) IN PATTERN_RULES:
        IF pattern matches cleaned OR description:
            RETURN result(label, confidence, stage="pattern")

    # Stage 2 -- kNN over multilingual MiniLM embeddings
    q <- embed(cleaned)                            # 384-dim vector
    sims <- cosine(q, train_embeddings)            # vs labelled set
    top_k <- argsort(sims)[-k:]                    # k = 5
    label <- majority_vote(train_labels[top_k])
    confidence <- mean(sims of winning label)
    IF confidence >= 0.65:
        RETURN result(label, confidence, stage="ml")

    # Stage 3 -- LLM fallback for the genuinely ambiguous remainder
    llm <- ask_local_llm_to_classify(description, cleaned)
    IF llm is valid:
        RETURN result(llm.label, llm.confidence, stage="llm")

    RETURN result(label, confidence, stage="ml_low_conf")
\end{lstlisting}

The ordering is the point. Stage one disposes of the unambiguous majority --- a
line containing ``WAZIRX'' is crypto, a line containing ``SALARY'' is salary ---
at essentially no cost. Stage two handles the noisy middle, where a multilingual
sentence embedding places ``Sharma Sweet Mart'' near other small-merchant
expenses even though no rule names it. Stage three, the only stage that invokes
the model, is reached by a small fraction of transactions, which keeps the
pipeline fast and keeps the model away from decisions a rule can make more
reliably. Every income-critical label --- capital market, crypto, interest,
dividend, freelance --- is one the rules and the embeddings handle with high
precision, so the costly stage is rarely the one deciding a taxable item.

\subsection{Reconciliation and the Unified Ledger}
With the bank credits classified, the Auditor extracts income items from all
three sources and merges them. Salary is matched across Form 16, the AIS, and
the bank; a difference below Rs.~100 is treated as confirmed, anything larger as
a mismatch. Every non-salary AIS entry is then checked against the bank and,
critically, against Form 16: an item the government reports but the employer's
Form 16 does not contain is exactly the kind of undeclared income that triggers
a notice, and it is recorded as an \texttt{ais\_only} item with a risk weight
attached. Bank credits that appear in neither Form 16 nor the AIS are recorded
too, but at a low weight, because the government does not yet know about them.

\subsection{Evidence-Weighted Risk Scoring}
The risk score is deliberately not a black box. It is an additive sum of
evidence-based weights, each chosen to reflect how strongly a given signal
predicts a notice, and then banded into four levels. The weights encode genuine
tax knowledge: undeclared crypto scores highest because the 194S TDS makes
non-disclosure almost certain to be caught; undeclared capital gains score
nearly as high because SEBI and STT data give the department full visibility; an
AIS-versus-Form-16 mismatch scores in the middle; a large cash deposit scores
lower.

\algname{compute\_risk(mismatches, anomalies) \quad with evidence WEIGHTS}
\begin{lstlisting}[style=pseudo]
WEIGHTS = { crypto_undeclared: 60, capital_gains_undeclared: 55,
            ais_mismatch_income: 25, freelance_undeclared: 25,
            property_unreported: 20, high_value_cash: 15,
            savings_interest_missing: 10 }

PROCEDURE compute_risk(mismatches, anomalies):
    total <- 0
    FOR m IN mismatches:
        total <- total + WEIGHTS[m.type]
    FOR a IN anomalies:
        IF a.in_ais:                  # already counted via mismatch
            CONTINUE                  # avoid double counting
        total <- total + a.risk_weight
    score <- min(100, total)
    IF   score < 20: level <- "LOW"
    ELIF score < 50: level <- "MEDIUM"
    ELIF score < 75: level <- "HIGH"
    ELSE:            level <- "CRITICAL"
    RETURN { score, level, breakdown }
\end{lstlisting}

The double-counting guard matters: an item that appears both as an AIS mismatch
and as a bank anomaly is counted once, not twice, so the score reflects distinct
pieces of evidence rather than the same fact seen from two angles. The banding
is intentionally conservative on the side of caution --- the design philosophy,
developed further in Chapter~\ref{chap:challenges}, is that a missed risk costs
the user money while an over-flag costs only a second look, so the thresholds
are set so that a genuine high-risk case is never quietly scored low.

\section{Module Two -- Optimizer Agent}
\label{sec:module-two}

The Optimizer answers the two questions a filer most wants answered: how much
tax, and under which regime. Its workflow, in Figure~\ref{fig:module2}, gathers
the applicable deductions, computes the liability under both regimes through the
CalculatorTool, compares them, and adds a CTC-restructuring analysis, finishing
with a natural-language explanation generated by the local model.

\begin{figure}[H]
    \centering
    \includegraphics[width=0.9\textwidth]{images/module2_flow.png}
    \caption{Optimizer Agent --- dual-regime computation via the deterministic engine, deduction and NPS optimisation, and regime recommendation.}
    \label{fig:module2}
\end{figure}

\subsection{The Deterministic Slab Engine}
The heart of the Optimizer --- and, by the ablation in
Chapter~\ref{chap:results}, the most important single component in the whole
system --- is the deterministic tax engine. It walks the slab structure for the
chosen regime, accumulating tax band by band. The same routine serves both
regimes; only the slab table, the standard deduction, and the rebate parameters
differ.

\algname{slab\_tax(taxable\_income, slabs) \quad / \quad new\_regime\_tax(gross, deductions)}
\begin{lstlisting}[style=pseudo]
PROCEDURE slab_tax(taxable_income, slabs):
    tax <- 0 ; remaining <- taxable_income
    FOR (lo, hi, rate) IN slabs:
        IF remaining <= 0: BREAK
        width <- (hi - lo + 1) IF hi != INF ELSE remaining
        in_band <- min(remaining, width)
        tax <- tax + in_band * rate
        remaining <- remaining - in_band
    RETURN round(tax, 2)

PROCEDURE new_regime_tax(gross, deductions):
    std <- min(75000, gross)             # FY25-26 new-regime std deduction
    allowed <- deductions.80CCD_2 + deductions.80CCH   # only these two
    taxable <- max(0, gross - std - allowed)
    base <- slab_tax(taxable, NEW_SLABS)

    # Section 87A rebate + marginal-relief zone
    rebate <- 0 ; relief <- FALSE
    IF taxable <= 1200000:
        rebate <- min(base, 60000)
    ELIF taxable <= 1275000:                 # marginal relief band
        excess <- taxable - 1200000
        IF base > excess:
            rebate <- base - excess ; relief <- TRUE

    after_rebate <- max(0, base - rebate)
    surcharge <- compute_surcharge(after_rebate, taxable, "new")
    cess <- 0.04 * (after_rebate + surcharge)
    RETURN after_rebate + surcharge + cess
\end{lstlisting}

Two subtleties are worth drawing out, because they are exactly the points at
which the general models in our benchmark fail. The first is the marginal-relief
band between Rs.~12,00,000 and roughly Rs.~12,75,000: inside it, the rebate is
not all-or-nothing but is tapered so that the tax never exceeds the income earned
above the twelve-lakh threshold. The second is the four percent cess, which is
applied to the sum of tax and surcharge, not to the base tax alone --- a step
the models routinely omit. The Old Regime routine mirrors this structure but
caps each Chapter VI-A deduction at its statutory limit (80C at Rs.~1,50,000,
24(b) at Rs.~2,00,000, and so on) before computing the slab tax against the
lower, deduction-friendly Old Regime slabs.

\subsection{HRA, Surcharge, and CTC Restructuring}
The HRA exemption under Section 10(13A) is computed as the statutory minimum of
three quantities, which is one of the places a naive implementation tends to
take only one of the three:

\algname{hra\_exemption(basic, hra\_received, rent\_paid, metro)}
\begin{lstlisting}[style=pseudo]
PROCEDURE hra_exemption(basic, hra_received, rent_paid, metro):
    IF rent_paid <= 0 OR hra_received <= 0: RETURN 0
    option1 <- hra_received
    option2 <- basic * (0.50 IF metro ELSE 0.40)
    option3 <- max(0, rent_paid - 0.10 * basic)
    RETURN min(option1, option2, option3)
\end{lstlisting}

The surcharge routine applies the appropriate rate band by total income and then
applies marginal relief so that the surcharge cannot push the total above the
income earned past the threshold. Finally, the CTC-restructuring analysis models
a concrete, lawful saving available even under the New Regime: an employer NPS
contribution under Section 80CCD(2), deductible up to ten percent of basic
salary. The Optimizer computes the tax with and without the maximum employer NPS
and reports the difference as an annual saving the user can ask their employer to
realise. The regime comparison itself is a straightforward selection of the
lower liability, but because both figures come from the deterministic engine, the
comparison is trustworthy in a way a model's self-reported numbers are not. The
language model is then asked only to phrase the recommendation, with an explicit
instruction not to compute anything; if the local model is unavailable, a
template fills in, so the numbers never depend on the model running at all.

\section{Module Three -- Compliance Agent}
\label{sec:module-three}

The Compliance agent decides which ITR form the filer must use and maps each
income item to its schedule. Its workflow, in Figure~\ref{fig:module3}, first
classifies the income types present, then applies the form-eligibility rules,
then maps schedules using document signals.

\begin{figure}[H]
    \centering
    \includegraphics[width=0.9\textwidth]{images/module3_flow.png}
    \caption{Compliance Agent --- ITR form selection and signal-based schedule mapping.}
    \label{fig:module3}
\end{figure}

\subsection{Form Selection as a Blocker Cascade}
ITR form selection is expressed not as a positive lookup but as a process of
elimination, which mirrors how the rules actually work: each form is disqualified
by the presence of certain income types, and the simplest surviving form is
chosen.

\algname{select\_itr\_form(income\_types)}
\begin{lstlisting}[style=pseudo]
PROCEDURE select_itr_form(income_types):
    # each form lists the income types that DISQUALIFY it
    blocked <- {}
    FOR form IN [ITR-1, ITR-2, ITR-4]:
        hits <- [t for t in RULES[form].blocked_by
                          if t in income_types]
        IF hits: blocked[form] <- hits

    IF "ITR-1" not in blocked AND
       ("salary" in income_types OR
        "other_sources" in income_types):
        RETURN "ITR-1"        # simplest case: salary/interest <= 50L
    ELIF "business_income" in income_types:
        RETURN ("ITR-4" if "ITR-4" not in blocked else "ITR-3")
    ELIF "ITR-2" not in blocked:
        RETURN "ITR-2"        # capital gains / crypto / foreign
    ELSE:
        RETURN "ITR-3"        # everything else
\end{lstlisting}

Capital gains or crypto, for instance, appear in the \texttt{blocked\_by} list
of both ITR-1 and ITR-4, so their presence forces the filer up to ITR-2 ---
which is precisely why the worked example earlier landed on ITR-2 the moment an
equity sale was detected.

\subsection{Signal-Based Schedule Inference}
The schedule mapping embodies one of the project's firmest design commitments:
a schedule is asserted only when a document actually supports it. This is most
visible in the handling of Schedule OS (income from other sources), which a less
careful system would attach to almost everyone. Here it is added only on a
concrete signal --- a bank interest credit, an AIS SFT-004 entry, an
old-regime Form 16 carrying deductions that imply interest-bearing investments,
or the presence of capital-gains or crypto activity that makes savings interest
an overwhelming empirical likelihood. The trade-off is explicit and intentional:
this keeps schedule precision very high at the cost of some recall, and the
reasoning is that asserting a schedule the documents do not support is a worse
error, for a filing tool, than omitting one the documents are silent about.
Where a broker or exchange CSV is supplied, the agent goes further and
auto-populates Schedule CG or Schedule VDA with computed gains, set-off, and the
194S TDS credit, drawing on the dedicated schedule builders.

\section{Module Four -- Critic Agent}
\label{sec:module-four}

The Critic is the system's conscience and the mechanism that most distinguishes
it from a prompted model. After the other three agents have run, it collects
every verifiable claim they produced, checks each one along several independent
axes, and either passes it or blocks it. Its workflow is shown in
Figure~\ref{fig:module4}.

\begin{figure}[H]
    \centering
    \includegraphics[width=0.9\textwidth]{images/module4_flow.png}
    \caption{Critic Agent --- claim verification against documents and the legal corpus, with the bounded re-run loop.}
    \label{fig:module4}
\end{figure}

\subsection{What Counts as a Claim, and How It Is Checked}
The Critic first assembles claims of several kinds: each deduction the Optimizer
applied (a ``deduction-eligibility'' claim), each legal section the narrative or
the schedule map references (a ``legal-reference'' claim), and each regime's
computed liability (an ``arithmetic'' claim). It then subjects each to up to four
checks.

\algname{verify\_claims(ctx)}
\begin{lstlisting}[style=pseudo]
PROCEDURE verify_claims(ctx):
    blocked <- [] ; verified <- []
    FOR claim IN collect_claims(ctx):

        # Check 1 -- regime compatibility (the classic hallucination)
        IF claim.type == "deduction_eligibility"
           AND ctx.recommended_regime == "new"
           AND normalize(claim.section) IN NEW_REGIME_BLOCKED:
            blocked.append(claim, reason="disallowed under New Regime")
            CONTINUE

        # Check 2 -- the cited section actually exists in the corpus
        IF NOT retriever.section_exists(claim.section):
            IF claim.type == "deduction_eligibility":
                blocked.append(claim, reason="fabricated section")
                CONTINUE

        # Check 3 -- arithmetic re-computation via the calculator
        IF claim.type == "arithmetic":
            recomputed <- calculator.recompute(claim, ctx)
            IF |recomputed - claim.value| > 100:
                warn("arithmetic_mismatch", claim)

        # Check 4 -- NLI faithfulness against retrieved legal text
        evidence <- retriever.retrieve(claim.section).text
        verdict  <- verifier.verify(claim.text, evidence)
        IF verdict.label == "HALLUCINATED":
            blocked.append(claim, reason="NLI contradiction")
        ELSE:
            verified.append(claim)

    status <- "success" IF blocked is empty
                         ELSE "needs_review"
    RETURN status, blocked, verified
\end{lstlisting}

The first check is the one that catches the textbook error --- an 80C, 80D, or
HRA deduction surviving into a New Regime recommendation --- by comparing the
claimed section against an explicit set of provisions the New Regime disallows.
The second guards against invented sections by requiring that the citation
resolve to a real node in the PageIndex corpus. The third independently
re-computes any stated tax figure through the same deterministic engine and
flags a divergence beyond Rs.~100. The fourth is the most general: it retrieves
the authoritative text of the cited section and asks an NLI cross-encoder
whether the claim is actually entailed by that text.

\subsection{The Faithfulness Verifier}
The verifier wraps a DeBERTa-based cross-encoder that scores a (premise,
hypothesis) pair --- here, (retrieved legal text, the claim) --- across the three
NLI labels, mapping entailment to \texttt{FAITHFUL}, contradiction to
\texttt{HALLUCINATED}, and neutral to \texttt{UNVERIFIED}. When the model cannot
be loaded, it degrades gracefully to a keyword-overlap heuristic rather than
failing, so verification still happens, if more crudely, on a minimal install.
The legal text it checks against comes from the PageIndex retriever, a tree-walk
over a hand-built corpus of more than eighty nodes covering both regimes, the
allowed and disallowed deductions, the capital-gains and crypto provisions, the
TDS sections, and the ITR-form rules; a section reference is matched directly to
its node by pattern, and a free-text query falls back to keyword scoring across
the tree. When the Critic blocks a claim, the orchestrator re-runs the
responsible agent --- most often the Optimizer, which then re-gathers deductions
with the blocked section excluded --- and the loop repeats, bounded at three
iterations, until the report is clean. As Chapter~\ref{chap:results} shows, this
single loop is what holds the hallucination rate down to a fraction of what the
same pipeline produces without it.
